\documentclass[12pt]{article}

\usepackage{amsmath}
\usepackage{parskip}
\usepackage{amssymb}
\usepackage{geometry}
\usepackage{hyperref}
\usepackage{booktabs}
\usepackage{hyperref}
\usepackage{tabularx}
\usepackage{float}
\usepackage{graphicx}
\usepackage{longtable}
\usepackage{url}
\usepackage{natbib}
\usepackage{pdfpages}
\usepackage{float}

\geometry{a4paper, margin=1in}

\title{Nichols Radial Injection Model (RIM) VII \\ The $\Lambda$CDM’s inconsistencies}
\author{Lawrence William Nichols}
\date{\today}

\begin{document}
\maketitle

\begin{abstract}
A compact catalog of only 22 spectroscopically characterized black holes, spanning stellar-mass to ultra-massive regimes, reveals universal chemical fingerprints (relativistic Fe K$\alpha$ lines, highly ionized winds, HCN/HCO$^+$ enhancements, warm Ne/O Goldilocks gas). When scaled by realistic sky-coverage ($\sim$1\% deep-field sampling) and detection-density factors ($\sim$22 named black holes vs. model inputs), the implied stellar-mass black hole population exceeds the entire baryonic mass budget of the observable universe under the standard $\Lambda$CDM 13.8 Gyr timeline. Adding progenitor stars and living stellar mass pushes the violation to 200--300\%. This mass-budget inconsistency falsifies $\Lambda$CDM's fixed cosmic age. The Nichols Radial Injection Model (RIM) resolves this via continuous radial mass injection from a static 4D bulk at an 11.07-year resonant pulse, extending effective cosmic time to 16.6 Gyr without exotic super-Eddington accretion or negative dark energy. RIM VII provides a mechanical, data-driven replacement for $\Lambda$CDM's timeline.
\end{abstract}

\section{Introduction}
The standard $\Lambda$CDM model faces mounting tension from JWST observations of mature early-universe structures and ultra-massive black holes (e.g., TON 618 at $z\approx2.2$, Phoenix A* cooling-flow anomalies). These require implausibly rapid growth within $\sim$500--900 Myr post-recombination.

\section{Prior Papers}
Prior RIM papers \citep{rimI,rimII,rimIII,rimIV,rimV,rimVI} establish a 16.6 Gyr hypersphere timeline with exogenous 11.07-year mass pulses ($\sim$10$^{44}$ kg/cycle) from a static 4D bulk, eliminating dark placeholders. Here we demonstrate that a short list of 22 black holes, scaled conservatively, forces a baryon-mass crisis incompatible with 13.8 Gyr.

\section{The 22-Object Spectroscopic Catalog}
The following table lists 22 well-studied black holes with diagnostic chemical/enrichment signatures:

\begin{table}[H]
\centering
\small
\caption{22 Black Hole Systems: Chemical Fingerprints}
\begin{tabular}{l l p{7.5cm}}
\toprule
Class & Object & Key Elements \& Data \\
\midrule
Stellar-mass & GRO J1655-40 ($\approx$6.3 M$_\odot$) & Fe, Ni, Cr, Ca; highly ionized wind; Fe XXV/XXVI absorption. \\ 
& Cygnus X-1 ($\approx$21.2 M$_\odot$) & H, He, Fe, O, N; Relativistic Fe K$\alpha$; super-solar Fe (2--5$\times$). \\
& LMC X-3 ($\approx$7 M$_\odot$) & H, He, C, N; low absorption; B-star dominated accretion. \\
& XTE J1118+480 ($\approx$7 M$_\odot$) & C, N, O; low metallicity; halo object with anomalous C/N ratios. \\
& MOA-2011-BLG-191 ($\approx$7.1 M$_\odot$) & H, He, Fe, O, N; Strong Relativistic Iron K$\alpha$ line \\
& V404 Cygni ($\approx$9 M$_\odot$) & H, He, Fe, Si, S, Mg; red-shifted outflows; neutral + ionized gas. \\
& GX 339-4 ($\approx$10 M$_\odot$) & Fe, Mg, Si; broad Fe K$\alpha$; high-spin disk dragging. \\
& XTE J1550-564 ($\approx$10 M$_\odot$) & Fe, Ca, S; Fe-K fluorescence; relativistic jet-disk flares. \\
& LMC X-1 ($\approx$10.9 M$_\odot$) & H, He, O, Ne, Si; persistent Fe K$\alpha$; hot plasma from O-star wind. \\

\midrule
IMBH & M82 X-1 ($\approx$400 M$_\odot$) & O, Ne, Fe, Mg; 6.4 keV Fe line; starburst molecular cloud reflection. \\
& NGC 5408 X-1 ($\approx$1,000 M$_\odot$) & He II, O, Fe; strong He II $\lambda4686$; extreme ionization. \\
& NGC 4395 ($\approx$10,000 M$_\odot$) & H, C, N, O, Ne, S, Fe; full AGN spectrum; “pocket quasar.” \\
& HLX-1 ($\approx$20,000 M$_\odot$) & H, He; H$\alpha$ emission; young massive cluster environment. \\
& NGC 404 ($\approx$500,000 M$_\odot$) & H$_2$, CO; cold molecular disk; ALMA CO(2--1) fueling. \\
\midrule
SMBH & Sagittarius A* ($\approx$4.3\,M M$_\odot$) & H, He, C, N, O, Ne, Mg, Si, Fe; HCN, CO, H$_2$O; Fe XXV/XXVI flares. \\
& NGC 1068 ($\approx$15\,M M$_\odot$) & HCN, HCO$^+$, SiO, S, N; ALMA HCN enhancement; obscured AGN tracer. \\
& CANUCS-LRD ($\approx$100\,M M$_\odot$) & H, C, N, h$_2$O; (low metallicity) High-velocity environment. \\
& NGC 1097 ($\approx$140\,M M$_\odot$) & HCN, HCO$^+$, CS; X-ray baked molecular gas. \\
& NGC 1275 ($\approx$800\,M M$_\odot$) & H$_2$, CO, Fe, Ne; warm filaments; jet-heated gas. \\
& M87* ($\approx$6.5\,B M$_\odot$) & O VIII, Fe, Mg, Ne, Si, S; enriched plasma; Fe $\approx$ solar. \\
& TON618 ($\approx$66\,B M$_\odot$) & Mg, Si, Fe, H, C, N; Optical/UV broad emission lines \\
& Phoenix A* ($\approx$100\,B M$_\odot$) & Ne, O, H; JWST warm Ne/O (300,000 K) Goldilocks gas. \\
\bottomrule
\end{tabular}
\label{tab:22objects}
\end{table}

These signatures (Fe reflection, ionized winds, X-ray-baked molecules) trace inner-galaxy metal factories from repeated star$\to$BH cycles.

\subsection{Scaling to the Unseen 99 \%: The Mass-Budget Crisis}

A compact catalog of only 22 spectroscopically well-characterized black holes, spanning stellar-mass to ultra-massive regimes (Table~1), already reveals universal chemical fingerprints (relativistic Fe K$\alpha$ lines, highly ionized winds, HCN/HCO$^+$ enhancements, warm Ne/O Goldilocks gas). When this sample is scaled by realistic deep-field sky-coverage ($\sim$1~\%) together with conservative detection-density and completeness factors, the implied total stellar-mass black-hole population reaches
\[
N_{\rm BH,total} = 4 \times 10^{21}.
\]

At an average mass of 10--20\,M$_\odot$ per object, the total mass locked in stellar-mass black holes becomes
\[
M_{\rm BH,total} \approx 4 \times 10^{22} - 8 \times 10^{22}\ M_\odot \tag{1}
\]

This already accounts for 70--160\,\% of the entire baryonic budget of the observable universe ($\sim$5--6 $\times 10^{22}$\,M$_\odot$, as inferred from CMB and Big-Bang nucleosynthesis constraints).  

Adding the mass in living stars ($\sim$5--8\,\% of the baryonic budget) plus the progenitor overhead required by a standard initial-mass-function (IMF) factor of 2--4$\times$ per black hole gives a processed-mass requirement of
\[
M_{\rm processed} \gtrsim (2 - 3) \times M_{\rm baryon}. \tag{2}
\]

The observable universe therefore contains 200--300\,\% more normal matter that has cycled through stars and black holes than is permitted by the standard baryon budget.  

To preserve the observed baryonic inventory without violation, the observable volume would have to scale by a factor of $\sim$200--300$\times$. In a standard expanding cosmology this implies a linear radius increase of $\sim$6$\times$ and a cosmic age of $\sim$80--100\,Gyr. The Nichols Radial Injection Model (RIM) instead resolves the tension exactly at an effective cosmic age of 16.6\,Gyr through continuous radial mass injection.

\section{Even the Full 95\% Dark Budget Cannot Rescue $\Lambda$CDM}

The standard $\Lambda$CDM model allocates only $\sim$5\% of the cosmic energy density to ordinary (baryonic) matter. The remaining $\sim$95\% is split between cold dark matter ($\sim$27\%) and dark energy ($\sim$68\%). It is often assumed these “dark variables” can absorb any mass-budget tension. They cannot.

Cold dark matter is collisionless by design. It cannot efficiently lose angular momentum, form dense accretion disks, or participate in the magnetohydrodynamic sorting and enrichment cycles required to produce the observed relativistic Fe K$\alpha$ lines, warm Ne/O Goldilocks gas (300,000 K), HCN/HCO$^+$ enhancements, or any of the chemical fingerprints listed in Table~\ref{tab:22objects}. Even generous assumptions about dark-matter accretion contribute negligibly to black-hole growth or heavy-element production.

Dark energy, modeled as a cosmological constant with negative pressure ($p = -\rho c^2$), behaves as perfectly uniform vacuum energy. In the exact Schwarzschild--de Sitter metric it does not flow toward or accrete into black holes. The equivalent mass contribution inside even a 100-billion-solar-mass event horizon is $\sim 10^{-40}\,M_\odot$ — utterly negligible and incapable of driving enrichment or resolving formation-time problems.

Thus, the entire mass-budget violation (200--300\% of available ordinary matter when the 22-object catalog is scaled realistically), the rapid chemical enrichment across all mass scales, and the timeline compression for ultra-massive objects (TON 618, Phoenix A*) fall squarely on the tiny 5\% baryonic component. The 95\% “dark” budget offers no physical mechanism to fix any of these problems.

This is not a minor tension. It is a structural failure of the entire ledger. The Nichols Radial Injection Model bypasses the impasse entirely by supplying fresh baryonic matter continuously from the static 4D bulk at the 11.07-year resonant pulse. No dark placeholders are required.

\section{Inconsistency with $\Lambda$CDM Timeline}
$\Lambda$CDM's 13.8 Gyr cannot accommodate:
\begin{itemize}
\item Early ultra-massive BHs (TON 618, Phoenix A*) without ad-hoc fixes.
\item Observed enrichment (fingerprints in Table~\ref{tab:22objects}) requiring multiple stellar generations.
\item Mass math above (independent of model assumptions).
\end{itemize}

\section{RIM Resolution}
RIM injects mass at 11.07 yr resonance \citep{rimV}:
\begin{equation}
M_u = N_{\rm pulses} \times 10^{44} \, \rm kg/pulse, \quad N_{\rm pulses} \approx 1.51 \times 10^9 \tag{3}
\end{equation}
The 16.6 Gyr timeline provides sufficient cycles for observed BH abundance/enrichment without violation.

\section{Conclusion and Predictions}
RIM VII falsifies $\Lambda$CDM's timeline via pure mass accounting from 22 objects. Predictions: Rubin/Euclid/ATHENA will reveal 200 billion M$_\odot$ candidates within 5 years, confirming extended time.

\section{Video Catalog}
A video overview of the full RIM series is available on the author's \href{https://www.youtube.com/@Nichols_Thought_Experiments}{Nichols Radial Injection Model (RIM)} YouTube channel.

\begin{thebibliography}{99}

\bibitem{rimI}
L.\,W.\,Nichols.
Nichols Radial Injection Model (RIM) and Radial ERB Inflows: A Mechanism for Progenitor-less Astrophysical Events, the Hubble Pulse, and 4D Substrate Interactions.
[Self-published], January 21, 2026.
\newline
\url{https://rxiverse.org/abs/2602.0036}

\bibitem{rimII}
L.\,W.\,Nichols.
Nichols Radial Injection Model (RIM) 11-Year Solar Cycle Pulses, GRB Milestones, Historical Alignments, and the 16.6 Gyr Hypersphere Universe.
[Self-published], February 2026.
\newline
\url{https://rxiverse.org/abs/2602.0038}

\bibitem{rimIII}
L.\,W.\,Nichols.
Nichols Radial Injection Model (RIM) III: The 3,000-Year Universal Odometer, the 1054 CE Primary Strike, and the Decaying 11-Year Cycle.
[Self-published], 15 February 2026.
\newline
\url{https://rxiverse.org/abs/2602.0043}

\bibitem{rimIV}
L.\,W.\,Nichols.
Nichols Radial Injection Model (RIM) IV: A 5D Mass-Regulated Manifold Solution to Cosmic Expansion.
[Self-published], 28 February 2026.
\newline
\url{https://rxiverse.org/abs/2603.0001}

\bibitem{rimV}
L.\,W.\,Nichols.
Nichols Radial Injection Model (RIM) V:
The Empirical Seal of the 11.07-Year Pulse
and Universal Resonance in Stellar Cycles
[Self-published], 4 March 2026.
\newline
\url{https://rxiverse.org/abs/2603.0003}

\bibitem{rimVI}
L.\,W.\,Nichols.
Nichols Radial Injection Model (RIM) VI:
The Birth of the Universe
[Self-published], 10 March 2026.
\newline
\url{https://rxiverse.org/abs/2603.0023}

\bibitem{einstein1935}
Einstein, A., \& Rosen, N. (1935).
\textit{The Particle Problem in the General Theory of Relativity}.
Physical Review, 48(1), 73.

\end{thebibliography}

\newpage

\section{Engineering Studies and Visual Data}
This section includes relevant visual diagrams.
\includepdf[pages=-]{Universe_Mass.pdf}
\end{document}